% Graph unlearning (GU) has emerged as a critical paradigm in the realm of graph-based machine unlearning, owing to its capability to remove sensitive information from trained graph neural networks (GNNs). Differing from traditional machine unlearning, GU approaches are meticulously crafted to untangle intricate interconnections scenarios, thereby revealing more meaningful perspectives on how to facilitate the targeted expunction of learned graph knowledge without compromising the integrity of model performance. Despite the proliferation of diverse GU strategies, the lack of standardized experimental setups and evaluation criteria has yielded inconsistencies in comparisons, hindering the thriving development of this research domain. To fill this gap, we have integrated 16 SOTA GU algorithms and 20 multi-domain datasets within OpenGU, enabling a variety of downstream tasks on 7 GNN backbones when responding to flexible unlearning requests. Based on this unified benchmark framework, we are able to provide a comprehensive and fair evaluation for the GU method, including effectiveness, efficiency, robustness, and scalability. Extensive experimental evaluation demonstrates the strong capability of GU in safeguarding graph data privacy while also uncovering the existing limitations, shedding light on potential avenues for future research.


Graph Machine Learning is essential for understanding and analyzing relational data. However, privacy-sensitive applications demand the ability to efficiently remove sensitive information from trained graph neural networks (GNNs), avoiding the unnecessary time and space overhead caused by retraining models from scratch.
% efficiently removing sensitive information from trained graph neural networks (GNNs), xxxx, has become a major challenge          
% without resorting to the costly and time-consuming process of retraining from scratch. 
To address this issue, Graph Unlearning (GU) has emerged as a critical solution, with the potential to support dynamic graph updates in data management systems and enable scalable unlearning in distributed data systems while ensuring privacy compliance.
Unlike machine unlearning in computer vision or other fields, GU faces unique difficulties due to the non-Euclidean nature of graph data and the recursive message-passing mechanism of GNNs. Additionally, the diversity of downstream tasks and the complexity of unlearning requests further amplify these challenges.
% Graph Machine Learning has become integral to understanding and analyzing complex relational data, highlighting the growing necessity for Graph Unlearning (GU) to address challenges in removing sensitive information from trained graph neural networks (GNNs). 
% 强调一下围绕图去做数据分析、处理（图机器学习）的重要性，然后顺理成章带出 graph unlearning
% Differing from machine unlearning in computer vision or other fields, GU presents unique challenges due to the intricate interactions within graph entities, the diversity of downstream tasks, and the varied nature of unlearning requests.
% 强调因为gu相比较于cv-based machine unlearning 或llm-based unlearning有独特的挑战（更加复杂、下游任务和unlearning request更加多样
Despite the proliferation of diverse GU strategies, the absence of a benchmark providing fair comparisons for GU, and the limited flexibility in combining downstream tasks and unlearning requests, have yielded inconsistencies in evaluations, hindering the development of this domain. To fill this gap, we present OpenGU, the first GU benchmark, where 16 SOTA GU algorithms and 37 multi-domain datasets are integrated, enabling various downstream tasks with 13 GNN backbones when responding to flexible unlearning requests. 
% 该领域在发展初期，许多方法没有在统一的基准下进行评估，因此需要开发一个同一的测试基准
 % opengu 都有什么
 Based on this unified benchmark framework, we are able to provide a comprehensive and fair evaluation for GU. Through extensive experimentation, we have drawn $8$ crucial conclusions about existing GU methods, while also gaining valuable insights into their limitations, shedding light on potential avenues for future research.
 % 做了很多实验得到了xxx结论得到了xxx insights等等